\documentclass[12pt,a4paper]{article}

% =========================================================
% HỖ TRỢ TIẾNG VIỆT
% =========================================================
\usepackage[utf8]{inputenc}
\usepackage[T5]{fontenc}
\usepackage[vietnamese]{babel}
\usepackage{lmodern}

% =========================================================
% GÓI TOÁN HỌC
% =========================================================
\usepackage{amsmath}
\usepackage{amssymb}
\usepackage{amsfonts}

% =========================================================
% VẼ ĐỒ THỊ VÀ HÌNH HỌC
% =========================================================
\usepackage{tikz}
\usepackage{pgfplots}

% Phiên bản tương thích của PGFPlots
\pgfplotsset{compat=1.18}

% Thư viện TikZ cần thiết
\usetikzlibrary{
    arrows.meta,
    intersections,
    calc,
    patterns,
    positioning
}

% =========================================================
% CĂN LỀ TRANG
% =========================================================
\usepackage[
    left=3cm,
    right=2cm,
    top=2.5cm,
    bottom=2.5cm
]{geometry}

% =========================================================
% ĐỊNH DẠNG ĐOẠN VĂN
% =========================================================
\usepackage{setspace}
\onehalfspacing

% =========================================================
% DANH SÁCH
% =========================================================
\usepackage{enumitem}

% =========================================================
% BẢNG
% =========================================================
\usepackage{array}
\usepackage{booktabs}
\usepackage{multirow}

% =========================================================
% HÌNH ẢNH
% =========================================================
\usepackage{graphicx}
\usepackage{float}

% =========================================================
% SIÊU LIÊN KẾT
% =========================================================
\usepackage[hidelinks]{hyperref}

% =========================================================
% BẮT ĐẦU VĂN BẢN
% =========================================================
\begin{document}


% =========================================================
% THÔNG TIN BÀI NGHIÊN CỨU
% =========================================================

\title{
    \textbf{ỨNG DỤNG HỆ BẤT PHƯƠNG TRÌNH HAI ẨN VÀO BÀI TOÁN TỐI ƯU HÓA CHI PHÍ THỰC ĐƠN DINH DƯỠNG}
}

\author{Thái Thị Mỹ Hoài}

\date{Tháng 9 năm 2026}

\maketitle
% =========================================================
% LỜI MỞ ĐẦU
% =========================================================

\section*{LỜI MỞ ĐẦU}

Quản lý chi tiêu cá nhân là một kỹ năng thiết yếu trong cuộc sống hiện nay. Đối với sinh viên sống xa nhà, việc cân đối ngân sách cho các nhu cầu học tập, sinh hoạt, đi lại là một bài toán khó. Trong đó, chi phí ăn uống thường chiếm tỷ trọng lớn và nếu không có kế hoạch hợp lý, sinh viên dễ rơi vào tình trạng chi tiêu vượt mức hoặc ăn uống thiếu dinh dưỡng. Do đó, việc xây dựng một thực đơn vừa tiết kiệm vừa đảm bảo dinh dưỡng cơ bản có ý nghĩa thực tiễn rất lớn.

Để giải quyết vấn đề này, toán học cung cấp công cụ đắc lực thông qua phương pháp mô hình hóa. Cụ thể, hệ bất phương trình hai ẩn giúp biểu diễn đồng thời các điều kiện ràng buộc như: giới hạn ngân sách, lượng calo, lượng protein tối thiểu và điều kiện số lượng thực phẩm không âm. Kết hợp với hàm mục tiêu biểu diễn tổng chi phí, bài toán thực tế sẽ trở thành bài toán tối ưu hóa.

Xuất phát từ ý tưởng trên, nhóm thực hiện đề tài:
\textbf{``Ứng dụng hệ bất phương trình hai ẩn vào bài toán tối ưu hóa chi phí thực đơn dinh dưỡng''}. 
Đề tài tập trung xây dựng mô hình đơn giản với hai biến quyết định, thiết lập hệ bất phương trình từ các số liệu giả định về giá thành và dinh dưỡng, từ đó xác định miền nghiệm để tìm ra phương án chi phí thấp nhất.

Nghiên cứu này không chỉ giúp người thực hiện củng cố kiến thức về hệ bất phương trình hai ẩn và phương pháp biểu diễn miền nghiệm, mà còn bước đầu hình thành tư duy tối ưu hóa. Đồng thời, đề tài khẳng định vai trò thực tế của toán học trong việc hỗ trợ lập kế hoạch chi tiêu và đưa ra các quyết định hợp lý trong cuộc sống hằng ngày.


% =========================================================
% CHƯƠNG 1
% =========================================================

\newpage

\section*{CHƯƠNG 1: GIỚI THIỆU ĐỀ TÀI}

% ---------------------------------------------------------
% 1.1
% ---------------------------------------------------------

\subsection*{1.1. Lý do chọn đề tài}

Sinh viên thường phải quản lý một khoản ngân sách có giới hạn trong suốt quá trình học tập. Trong đó, ăn uống là nhu cầu thiết yếu và chiếm một phần đáng kể trong tổng chi tiêu. Tuy nhiên, nếu lựa chọn thực phẩm không hợp lý, sinh viên có thể gặp hai vấn đề: hoặc chi phí ăn uống quá cao, hoặc thực đơn có giá thành thấp nhưng không đáp ứng được những yêu cầu dinh dưỡng cơ bản.
Từ thực tế đó, một vấn đề được đặt ra là:
\begin{center}
\textit{``Làm thế nào để xây dựng một thực đơn có chi phí thấp nhất nhưng vẫn đáp ứng được các yêu cầu tối thiểu về dinh dưỡng và phù hợp với ngân sách của một sinh viên?''}
\end{center}
Đây là một bài toán có thể được mô hình hóa bằng hệ bất phương trình hai ẩn. Nếu lựa chọn hai loại thực phẩm hoặc hai nhóm thực phẩm làm cơ sở nghiên cứu, ta có thể sử dụng hai biến $x$ và $y$ để biểu diễn số lượng của chúng. Các điều kiện về chi phí, năng lượng, protein và số lượng thực phẩm có thể được chuyển thành các bất phương trình tuyến tính theo $x$ và $y$.

Sau khi xây dựng hệ bất phương trình, miền nghiệm có thể được biểu diễn trực quan trên mặt phẳng tọa độ $Oxy$. Kết hợp với hàm chi phí, ta có thể xác định phương án phù hợp và tìm phương án có tổng chi phí nhỏ nhất.

Việc lựa chọn đề tài này không chỉ giúp củng cố kiến thức về hệ bất phương trình hai ẩn mà còn cho thấy khả năng ứng dụng của toán học vào một vấn đề gần gũi với đời sống sinh viên. Thông qua quá trình xây dựng và giải mô hình, người thực hiện có cơ hội rèn luyện tư duy logic, khả năng phân tích điều kiện và kỹ năng đưa ra quyết định dựa trên cơ sở toán học.


% ---------------------------------------------------------
% 1.2
% ---------------------------------------------------------

\subsection*{1.2. Mục tiêu của đề tài}

Đề tài được thực hiện với các mục tiêu chính sau:
\begin{itemize}
    \item Tìm hiểu cơ sở lý thuyết về hệ bất phương trình bậc nhất hai ẩn và bài toán tối ưu hóa.
    \item Tìm hiểu cách biểu diễn một hệ bất phương trình trên mặt phẳng tọa độ $Oxy$.
    \item Xây dựng mô hình toán học cho bài toán tối ưu hóa chi phí thực đơn của một sinh viên.
    \item Xác định các biến quyết định đại diện cho số lượng thực phẩm được sử dụng trong thực đơn.
    \item Chuyển các điều kiện về chi phí và dinh dưỡng thành một hệ bất phương trình hai ẩn.
    \item Thiết lập hàm mục tiêu biểu diễn tổng chi phí thực đơn cần tối thiểu hóa.
    \item Xác định miền nghiệm khả thi của hệ bất phương trình.
    \item Sử dụng phương pháp đồ thị và các phép tính cần thiết để tìm phương án thực đơn có chi phí thấp nhất.
    \item Phân tích ý nghĩa kinh tế và dinh dưỡng của nghiệm tối ưu.
    \item Đánh giá ưu điểm, hạn chế và khả năng thay đổi của mô hình khi các điều kiện thực tế thay đổi.
\end{itemize}


% ---------------------------------------------------------
% 1.3
% ---------------------------------------------------------

\subsection*{1.3. Đối tượng nghiên cứu}

Đối tượng nghiên cứu của đề tài là việc ứng dụng hệ bất phương trình hai ẩn vào bài toán tối ưu hóa chi phí thực đơn cho một sinh viên.

Trong mô hình toán học, hai biến $x$ và $y$ được sử dụng để biểu diễn số lượng của hai loại thực phẩm hoặc hai nhóm thực phẩm được lựa chọn trong thực đơn.

Các yếu tố chính được xem xét trong mô hình gồm:

\begin{itemize}
    \item Tổng chi phí thực phẩm.
    \item Lượng năng lượng cung cấp.
    \item Lượng protein cung cấp.
    \item Các giới hạn về số lượng thực phẩm.
    \item Điều kiện ngân sách của sinh viên.
    \item Điều kiện không âm của các biến quyết định.
\end{itemize}
Từ các yếu tố trên, bài toán được xây dựng thành một hệ bất phương trình hai ẩn kết hợp với hàm mục tiêu chi phí.


% ---------------------------------------------------------
% 1.4
% ---------------------------------------------------------

\subsection*{1.4. Phạm vi nghiên cứu}

Đề tài được giới hạn trong việc xây dựng một mô hình toán học đơn giản nhằm minh họa bài toán tối ưu hóa chi phí thực đơn cho một sinh viên.

Trong mô hình, đề tài tập trung nghiên cứu hai biến quyết định tương ứng với hai loại thực phẩm hoặc hai nhóm thực phẩm giả định. Các tiêu
chí chính được xem xét gồm chi phí, năng lượng và protein. Tùy theo mô hình cụ thể, có thể bổ sung thêm các điều kiện về carbohydrate, chất béo hoặc số lượng suất ăn.

Các số liệu về giá thành và thành phần dinh dưỡng được giả định nhằm phục vụ mục đích học tập. Các số liệu này không nhằm đưa ra khuyến nghị dinh dưỡng thực tế cho mọi sinh viên.

Đề tài chưa xét đến một số yếu tố phức tạp trong thực tế như:

\begin{itemize}
    \item Sự thay đổi giá thực phẩm theo thời gian.
    \item Sở thích và khẩu vị của từng sinh viên.
    \item Dị ứng hoặc yêu cầu dinh dưỡng riêng của từng người.
    \item Vitamin, khoáng chất và các vi chất dinh dưỡng.
    \item Sự thay đổi thành phần dinh dưỡng trong quá trình chế biến.
    \item Hao hụt thực phẩm khi nấu ăn.
    \item Sự khác nhau về nhu cầu năng lượng giữa từng cá nhân.
\end{itemize}

Việc giới hạn phạm vi nghiên cứu giúp mô hình có cấu trúc đơn giản, phù hợp với việc giải bằng hệ bất phương trình hai ẩn và phương pháp đồ thị.


% ---------------------------------------------------------
% 1.5
% ---------------------------------------------------------

\subsection*{1.5. Phương pháp nghiên cứu}

Đề tài sử dụng chủ yếu phương pháp mô hình hóa toán học, kết hợp với phương pháp giải hệ bất phương trình và phương pháp đồ thị.


% ---------------------------------------------------------
% 1.6
% ---------------------------------------------------------

\subsection*{1.6. Ý nghĩa của đề tài}

Trước hết, đề tài mang lại ý nghĩa thiết thực về mặt học tập và toán học. Thay vì chỉ tiếp cận lý thuyết thuần túy, việc vận dụng hệ bất phương trình hai ẩn vào một tình huống gần gũi với đời sống sinh viên giúp người thực hiện hiểu rõ hơn mối quan hệ giữa các biến, các điều kiện ràng buộc và hàm mục tiêu. Đồng thời, việc biểu diễn miền nghiệm trên mặt phẳng tọa độ Oxy cũng giúp trực quan hóa quá trình giải bài toán tối ưu.

Về mặt thực tiễn, mô hình cho thấy vai trò của toán học trong việc hỗ trợ lập kế hoạch chi tiêu. Khi ngân sách có giới hạn nhưng vẫn phải đáp ứng các tiêu chuẩn nhất định, mô hình toán học sẽ cung cấp cơ sở khoa học để đưa ra quyết định phù hợp thay vì lựa chọn theo cảm tính.

Tuy nhiên, do sử dụng các số liệu giả định, kết quả của mô hình chủ yếu mang tính chất minh họa cho phương pháp tối ưu hóa chi phí. Trong thực tế, việc xây dựng một thực đơn hoàn chỉnh sẽ đòi hỏi phải xem xét nhiều yếu tố dinh dưỡng và nhu cầu cá nhân phức tạp hơn.

Nhìn chung, đề tài đã góp phần kết nối kiến thức toán học với một vấn đề thực tế quen thuộc, qua đó rèn luyện khả năng xây dựng mô hình, phân tích dữ liệu và tìm kiếm phương án tối ưu dựa trên các điều kiện đã cho.
% =========================================================
% CHƯƠNG 2: CƠ SỞ LÝ LUẬN
% =========================================================

\newpage

\section*{CHƯƠNG 2: CƠ SỞ LÝ LUẬN}

\subsection*{2.1. Bất phương trình bậc nhất hai ẩn}

Bất phương trình bậc nhất hai ẩn là bất phương trình có dạng:
\[
ax+by\leq c
\]
hoặc:
\[
ax+by\geq c,
\]
trong đó $a,b,c$ là các số thực và $a,b$ không đồng thời bằng $0$.

Cặp số $(x;y)$ được gọi là nghiệm của bất phương trình nếu khi thay
$x$ và $y$ bằng các giá trị tương ứng vào bất phương trình, ta nhận được một mệnh đề đúng.

Trong các bài toán thực tế, bất phương trình bậc nhất hai ẩn được sử dụng để biểu diễn những điều kiện giới hạn của hai đại lượng.

\subsection*{2.2. Hệ bất phương trình bậc nhất hai ẩn}

Hệ bất phương trình bậc nhất hai ẩn là tập hợp gồm nhiều bất phương trình cùng chứa hai ẩn $x$ và $y$.

Một hệ bất phương trình có thể có dạng:
\[
\left\{
\begin{array}{l}
a_1x+b_1y\leq c_1,\\
a_2x+b_2y\geq c_2,\\
x\geq0,\quad y\geq0.
\end{array}
\right.
\]

Nghiệm của hệ là những cặp số $(x;y)$ đồng thời thỏa mãn tất cả các bất phương trình trong hệ.

Trong bài toán thực tế, hệ bất phương trình giúp biểu diễn nhiều điều kiện khác nhau cùng một lúc.

\subsection*{2.3. Miền nghiệm của hệ bất phương trình}

Mỗi bất phương trình bậc nhất hai ẩn xác định một miền nghiệm trên mặt phẳng tọa độ $Oxy$.

Đường thẳng:

\[
ax+by=c
\]

được gọi là đường biên của bất phương trình. Đường thẳng này chia mặt phẳng thành hai nửa mặt phẳng, trong đó một nửa là miền nghiệm của bất phương trình.

Miền nghiệm của một hệ bất phương trình là phần giao của các miền nghiệm tương ứng.

Đối với bài toán thực đơn, mỗi điểm thuộc miền nghiệm có thể được hiểu là một phương án lựa chọn thực phẩm thỏa mãn các điều kiện đã đặt ra.

\subsection*{2.4. Hàm mục tiêu và bài toán tối ưu}

Trong một bài toán tối ưu, cần xác định một đại lượng cần tìm giá trị lớn nhất hoặc nhỏ nhất. Đại lượng đó được biểu diễn bởi hàm mục tiêu.

Đối với bài toán tối ưu hóa chi phí thực đơn, hàm mục tiêu là tổng chi phí mua thực phẩm và mục tiêu là tìm giá trị nhỏ nhất:
\[
Z\longrightarrow\min.
\]
Các điều kiện về số lượng thực phẩm, ngân sách và dinh dưỡng được sử dụng làm các ràng buộc của bài toán.

Như vậy, bài toán tối ưu có thể được hiểu đơn giản là:
\[
\boxed{
\text{Tìm phương án phù hợp với các điều kiện và có chi phí nhỏ nhất.}
}
\]

\subsection*{2.5. Một số điều kiện trong bài toán thực đơn}

Khi xây dựng bài toán thực đơn dinh dưỡng, có thể xem xét một số điều kiện cơ bản như:
\begin{itemize}
    \item Chi phí thực đơn không vượt quá ngân sách cho phép.
    \item Lượng năng lượng cung cấp đạt mức tối thiểu.
    \item Lượng protein đạt mức yêu cầu.
    \item Lượng thực phẩm không được âm.
\end{itemize}

Các điều kiện trên có thể được biểu diễn bằng các bất phương trình.

Việc kết hợp các điều kiện này tạo thành một hệ bất phương trình, từ đó xác định các phương án thực đơn có thể chấp nhận.

\subsection*{2.6. Ý nghĩa của việc áp dụng toán học vào bài toán thực đơn}

Việc sử dụng hệ bất phương trình hai ẩn giúp chuyển một vấn đề thực tế thành một bài toán toán học cụ thể.

Thông qua mô hình toán học, có thể xác định được những phương án đáp ứng các yêu cầu đặt ra và lựa chọn phương án phù hợp nhất theo mục tiêu của bài toán.

Trong đề tài này, kiến thức về hệ bất phương trình hai ẩn được sử dụng làm cơ sở để xây dựng và giải bài toán tối ưu hóa chi phí thực đơn cho sinh viên.


% =========================================================
% KẾT THÚC CHƯƠNG 2
% =========================================================
% =========================================================
% CHƯƠNG 3
% =========================================================

\newpage

\section*{CHƯƠNG 3: XÂY DỰNG MÔ HÌNH VÀ THIẾT LẬP HỆ BẤT PHƯƠNG TRÌNH}

\subsection*{3.1. Tình huống thực tế}

Trong thực tế, sinh viên thường phải cân đối giữa nhu cầu ăn uống, dinh dưỡng và khả năng tài chính. Đặc biệt đối với sinh viên sống xa gia đình, việc lập kế hoạch chi tiêu cho các bữa ăn trong tuần có thể giúp sử dụng ngân sách hợp lý hơn.

Tuy nhiên, một sinh viên không nhất thiết phải sử dụng duy nhất một loại thực đơn trong suốt một tuần. Để mô hình gần với thực tế hơn, đề tài giả định sinh viên có thể lựa chọn giữa hai loại suất ăn đại diện:
\begin{itemize}
    \item \textbf{Thực đơn A:} suất cơm trứng và rau xanh, có chi phí thấp hơn và cung cấp mức năng lượng, protein và rau xanh ở mức cơ bản.
    \item \textbf{Thực đơn B:} suất cơm gà và rau xanh, có chi phí cao hơn nhưng cung cấp nhiều năng lượng, protein và rau xanh hơn.
\end{itemize}

Hai loại thực đơn trên chỉ mang tính đại diện cho hai mức lựa chọn về chi phí và dinh dưỡng. Trong thực tế, sinh viên có thể thay đổi món ăn hằng ngày. Mô hình được xây dựng nhằm minh họa cách sử dụng hệ bất phương trình hai ẩn để lựa chọn số lượng suất ăn sao cho vừa đáp ứng nhu cầu dinh dưỡng, vừa kiểm soát chi phí.

Để phù hợp với phạm vi của đề tài, các số liệu sử dụng trong mô hình là số liệu giả định phục vụ mục đích học tập, không phải khuyến nghị dinh dưỡng áp dụng cho mọi sinh viên.

\subsection*{3.2. Các giả định của mô hình}

Để xây dựng mô hình toán học đơn giản và có thể giải bằng phương pháp đồ thị, đề tài sử dụng các giả định sau:
\begin{enumerate}
    \item Thời gian xem xét là một tuần, tương ứng với 7 ngày.
    \item Sinh viên có thể lựa chọn hai loại suất ăn đại diện là thực đơn A và thực đơn B.
    \item Mỗi suất ăn cung cấp một lượng năng lượng, protein và rau xanh tương đối ổn định theo bảng số liệu giả định.
    \item Chi phí của mỗi suất ăn được xem là cố định trong phạm vi thời gian nghiên cứu.
    \item Tổng nhu cầu dinh dưỡng được xét trên toàn bộ một tuần thay vì yêu cầu mỗi ngày phải sử dụng đúng một loại thực đơn.
    \item Mô hình chỉ xét một số yếu tố cơ bản gồm chi phí, năng lượng, protein và rau xanh. Các yếu tố như vitamin cụ thể, khoáng chất, chất béo, khẩu vị và sự thay đổi giá thực phẩm chưa được đưa vào.
    \item Các số liệu trong mô hình có thể được điều chỉnh khi ngân sách hoặc nhu cầu dinh dưỡng thay đổi. Đây là cơ sở để đánh giá  tính linh hoạt của mô hình.
\end{enumerate}

Các giả định trên giúp mô hình vừa đơn giản, phù hợp với kiến thức về hệ bất phương trình hai ẩn, vừa có khả năng mô phỏng một tình huống thực tế trong quản lý chi phí ăn uống của sinh viên.

\subsection*{3.3. Bảng dữ liệu của mô hình}

Dựa trên các giả định đã xây dựng, bảng dữ liệu của hai loại suất ăn được trình bày như sau:

\begin{center}
\textbf{Bảng 3.1. Chi phí và giá trị dinh dưỡng của hai loại suất ăn}
\end{center}

\begin{center}
\renewcommand{\arraystretch}{1.4}

\begin{tabular}{|p{4.0cm}|c|c|c|}
\hline
\textbf{Tiêu chí}
&
\textbf{Thực đơn A}
&
\textbf{Thực đơn B}
&
\textbf{Nhu cầu trong tuần}
\\
\hline

Chi phí (VNĐ/suất)
&
25\,000
&
35\,000
&
Không vượt quá 600\,000
\\
\hline

Năng lượng (kcal/suất)
&
600
&
700
&
Ít nhất 10\,500
\\
\hline

Protein (g/suất)
&
20
&
35
&
Ít nhất 420
\\
\hline

Rau xanh (g/suất)
&
100
&
150
&
Ít nhất 2\,100
\\
\hline

\end{tabular}
\end{center}

Các giá trị trong bảng được xây dựng nhằm tạo ra một bài toán có thể giải bằng hệ bất phương trình hai ẩn. Trong thực tế, các giá trị này có thể được thay thế bằng dữ liệu dinh dưỡng và giá thành cụ thể của các món ăn mà sinh viên lựa chọn.

\subsection*{3.4. Xác định biến quyết định}

Gọi:
\[
x = \text{số suất thực đơn A được sử dụng trong tuần},
\]
\[
y = \text{số suất thực đơn B được sử dụng trong tuần}.
\]
Do $x$ và $y$ biểu diễn số lượng suất ăn nên:
\[
x\geq0,\qquad y\geq0.
\]
Nếu xét số suất ăn thực tế, $x$ và $y$ có thể được xem là các số nguyên không âm:
\[
x,y\in\mathbb{N}.
\]
Tuy nhiên, khi giải bài toán bằng phương pháp đồ thị, ta có thể trước hết xem $x$ và $y$ là các biến thực không âm để xác định miền nghiệm liên tục. Sau đó có thể lựa chọn các điểm nguyên gần nghiệm tối ưu để đưa ra phương án thực tế.

Cách xây dựng này giúp mô hình vừa thuận tiện cho việc biểu diễn hình học, vừa có thể chuyển thành phương án số suất ăn cụ thể.

\subsection*{3.5. Thiết lập các ràng buộc của mô hình}

\subsubsection*{3.5.1. Ràng buộc về ngân sách}

Sinh viên có ngân sách tối đa là $600\,000$ VNĐ cho các suất ăn được xét trong một tuần.

Chi phí của $x$ suất thực đơn A là: \textbf{25\,000x.}

Chi phí của $y$ suất thực đơn B là: \textbf{35\,000y.}

Do tổng chi phí không được vượt quá ngân sách:
\[
25\,000x+35\,000y\leq600\,000.
\]
Chia hai vế cho $5\,000$, ta được:
\[
\boxed{
5x+7y\leq120
}
\tag{1}
\]

\subsubsection*{3.5.2. Ràng buộc về năng lượng}

Mỗi suất thực đơn A cung cấp $600$ kcal và mỗi suất thực đơn B cung
cấp $700$ kcal.

Tổng năng lượng nhận được trong tuần là:
\[
600x+700y.
\]
Để đạt mức năng lượng tối thiểu giả định là $10\,500$ kcal:
\[
600x+700y\geq10\,500.
\]
Chia hai vế cho $100$, ta được:
\[
\boxed{
6x+7y\geq105
}
\tag{2}
\]

\subsubsection*{3.5.3. Ràng buộc về protein}

Mỗi suất thực đơn A cung cấp $20$ g protein và mỗi suất thực đơn B cung cấp $35$ g protein.

Tổng lượng protein là:
\[
20x+35y.
\]
Theo mức tối thiểu giả định:
\[
20x+35y\geq420.
\]
Chia hai vế cho $5$, ta được:
\[
\boxed{
4x+7y\geq84
}
\tag{3}
\]

\subsubsection*{3.5.4. Ràng buộc về rau xanh}

Mỗi suất thực đơn A có $100$ g rau xanh và mỗi suất thực đơn B có $150$ g rau xanh.

Tổng lượng rau xanh trong tuần là:
\[
100x+150y.
\]
Yêu cầu tối thiểu là $2\,100$ g:
\[
100x+150y\geq2\,100.
\]
Chia hai vế cho $50$, ta được:
\[
\boxed{
2x+3y\geq42
}
\tag{4}
\]
Ràng buộc này được sử dụng như một đại diện đơn giản cho yêu cầu bổ sung rau xanh trong khẩu phần. Trong thực tế, nhu cầu dinh dưỡng cần được xem xét đầy đủ hơn, bao gồm nhiều nhóm thực phẩm và chất dinh dưỡng khác nhau.

\subsubsection*{3.5.5. Điều kiện không âm}

Vì số lượng suất ăn không thể âm nên:
\[
\boxed{
x\geq0,\qquad y\geq0
}
\tag{5}
\]

\subsection*{3.6. Hệ bất phương trình của mô hình}

Từ các ràng buộc trên, bài toán được mô hình hóa bởi hệ bất phương trình:

\[
\boxed{
\left\{
\begin{array}{ll}
5x+7y\leq120, & \text{(Ngân sách)},\\[4pt]
6x+7y\geq105, & \text{(Năng lượng)},\\[4pt]
4x+7y\geq84, & \text{(Protein)},\\[4pt]
2x+3y\geq42, & \text{(Rau xanh)},\\[4pt]
x\geq0,\quad y\geq0. & \text{(Không âm)}
\end{array}
\right.
}
\]

Tập hợp tất cả các điểm $(x,y)$ thỏa mãn đồng thời các bất phương trình trên tạo thành miền nghiệm của bài toán trên mặt phẳng tọa độ
$Oxy$.

Mỗi điểm $(x,y)$ thuộc miền nghiệm có thể được hiểu là một phương án lựa chọn số lượng suất ăn A và B thỏa mãn đồng thời các điều kiện về ngân sách và dinh dưỡng.

\subsection*{3.7. Xây dựng hàm mục tiêu chi phí}

Mục tiêu của đề tài là tìm phương án có tổng chi phí thấp nhất nhưng vẫn đáp ứng các yêu cầu đã đặt ra.

Tổng chi phí của $x$ suất thực đơn A và $y$ suất thực đơn B là:
\[
C=25\,000x+35\,000y.
\]
Do cần tối thiểu hóa chi phí, hàm mục tiêu được viết:
\[
\boxed{
\min C=25\,000x+35\,000y
}
\]
Như vậy, bài toán tối ưu của đề tài là tìm điểm $(x,y)$ thuộc miền nghiệm của hệ bất phương trình sao cho giá trị của hàm $C$ đạt nhỏ nhất.

\subsection*{3.8. Kiểm tra tính khả thi của mô hình}

Trước khi tiến hành tối ưu hóa, cần kiểm tra xem hệ bất phương trình có tồn tại nghiệm hay không.

Xét phương án:
\[
(x,y)=(9,8).
\]
Khi đó:
\[
5(9)+7(8)=45+56=101\leq120,
\]
\[
6(9)+7(8)=54+56=110\geq105,
\]
\[
4(9)+7(8)=36+56=92\geq84,
\]
\[
2(9)+3(8)=18+24=42.
\]
Đồng thời:
\[
9\geq0,\qquad8\geq0.
\]
Như vậy, điểm $(9,8)$ là một nghiệm khả thi của hệ bất phương trình.

Chi phí tương ứng là:
\[
C=25\,000(9)+35\,000(8)
=225\,000+280\,000
=505\,000\ \text{VNĐ}.
\]
Điều này chứng tỏ miền nghiệm của bài toán không rỗng và bài toán có thể tiếp tục được giải bằng phương pháp đồ thị hoặc công cụ tính toán.

\subsection*{3.9. Tính linh hoạt của mô hình}

Một ưu điểm của mô hình là các tham số có thể được điều chỉnh khi điều kiện thực tế thay đổi.

Chẳng hạn, nếu ngân sách của sinh viên thay đổi từ $600\,000$ VNĐ lên
$700\,000$ VNĐ thì ràng buộc ngân sách trở thành:
\[
25\,000x+35\,000y\leq700\,000.
\]
Nếu nhu cầu năng lượng thay đổi, hệ số ở vế phải của bất phương trình năng lượng cũng được thay đổi tương ứng.

Ví dụ, nếu mức năng lượng tối thiểu được điều chỉnh thành
$11\,000$ kcal thì:
\[
600x+700y\geq11\,000.
\]
Khi đó miền nghiệm của bài toán có thể thay đổi và phương án tối ưu cũng có khả năng thay đổi.

Ngoài ra, khi giá thực phẩm thay đổi, hàm mục tiêu cũng thay đổi. Ví dụ, nếu giá thực đơn A tăng lên $27\,000$ VNĐ/suất thì hàm mục tiêu trở thành:
\[
C=27\,000x+35\,000y.
\]
Như vậy, cùng một hệ ràng buộc dinh dưỡng nhưng giá thực phẩm khác nhau có thể dẫn đến phương án lựa chọn khác nhau.

Điều này cho thấy mô hình có tính linh hoạt và có thể được điều chỉnh theo điều kiện tài chính hoặc nhu cầu dinh dưỡng của sinh viên.

\subsection*{3.10. Ý nghĩa của mô hình}

Mô hình được xây dựng cho phép chuyển một tình huống quản lý chi tiêu ăn uống của sinh viên thành một bài toán toán học cụ thể.

Hai biến $x$ và $y$ đại diện cho số lượng hai loại suất ăn. Các yêu cầu về ngân sách, năng lượng, protein và rau xanh được biểu diễn bằng các bất phương trình tuyến tính. Trong khi đó, tổng chi phí được biểu diễn bằng một hàm mục tiêu cần tối thiểu hóa.

Nhờ đó, bài toán thực tế:
\begin{center}
``Lựa chọn thực đơn phù hợp với ngân sách và nhu cầu dinh dưỡng''
\end{center}
được chuyển thành bài toán:
\begin{center}
``Tìm giá trị nhỏ nhất của một hàm tuyến tính trên miền nghiệm của một hệ bất phương trình hai ẩn''.
\end{center}
Mô hình cũng cho phép quan sát trực quan sự thay đổi của miền nghiệm khi ngân sách, giá thực phẩm hoặc nhu cầu dinh dưỡng thay đổi.

\subsection*{3.11. Hạn chế của mô hình}

Mặc dù có tính minh họa và phù hợp với mục tiêu nghiên cứu của đề tài, mô hình vẫn có một số hạn chế.

Thứ nhất, mô hình chỉ sử dụng hai loại suất ăn trong khi thực tế sinh viên có thể lựa chọn nhiều loại thực phẩm và món ăn khác nhau.

Thứ hai, giá trị dinh dưỡng của mỗi suất ăn được xem là cố định, trong khi thành phần dinh dưỡng thực tế có thể thay đổi tùy thuộc vào nguyên liệu, khẩu phần và cách chế biến.

Thứ ba, mô hình chỉ xét một số yếu tố dinh dưỡng cơ bản như năng lượng, protein và rau xanh. Các chất dinh dưỡng khác như chất béo, vitamin và khoáng chất chưa được đưa vào.

Thứ tư, mô hình chưa xét đến sở thích cá nhân, sự đa dạng của bữa ăn và các yếu tố khác ảnh hưởng đến quyết định lựa chọn thực phẩm.

Do đó, kết quả của mô hình chủ yếu có ý nghĩa minh họa cho việc ứng dụng hệ bất phương trình hai ẩn và tư duy tối ưu hóa vào một tình huống thực tế.

\subsection*{3.12. Kết luận chương}

Trong Chương 3, đề tài đã xây dựng mô hình toán học cho bài toán tối ưu hóa chi phí thực đơn dinh dưỡng dành cho sinh viên.

Các biến quyết định $x$ và $y$ được sử dụng để biểu diễn số lượng hai loại suất ăn. Từ các giả định về giá thành và giá trị dinh dưỡng, các điều kiện về ngân sách, năng lượng, protein và rau xanh được chuyển thành hệ bất phương trình hai ẩn:
\[
\left\{
\begin{array}{l}
5x+7y\leq120,\\
6x+7y\geq105,\\
4x+7y\geq84,\\
2x+3y\geq42,\\
x\geq0,\quad y\geq0.
\end{array}
\right.
\]
Hàm mục tiêu của bài toán là:
\[
\boxed{
\min C=25\,000x+35\,000y.
}
\]
Mô hình có nghiệm khả thi và có thể được biểu diễn trên mặt phẳng tọa độ $Oxy$. Đây là cơ sở để sang chương tiếp theo tiến hành biểu diễn đồ thị, xác định miền nghiệm, tìm các điểm cực biên và xác định phương án có chi phí tối ưu.

Đồng thời, mô hình có tính linh hoạt khi các tham số như ngân sách, giá thực phẩm hoặc nhu cầu dinh dưỡng thay đổi. Điều này giúp mô hình không chỉ phục vụ mục đích minh họa kiến thức toán học mà còn thể hiện được khả năng ứng dụng của hệ bất phương trình hai ẩn vào bài toán quản lý chi tiêu trong đời sống.

% =========================================================
% KẾT THÚC CHƯƠNG 3
% =========================================================
% =========================================================
% CHƯƠNG 4
% =========================================================

\newpage

\section*{CHƯƠNG 4: GIẢI VÀ PHÂN TÍCH BÀI TOÁN}

\subsection*{4.1. Bài toán}

Từ mô hình đã xây dựng ở Chương 3, gọi $x$ là số suất Thực đơn A và $y$ là số suất Thực đơn B được sử dụng trong tuần.

Hệ bất phương trình của bài toán là:

\[
\begin{cases}
5x+7y\leq120,\\
5x+7y\geq100,\\
4x+7y\geq90,\\
x\geq0,\\
y\geq0.
\end{cases}
\]

Hàm mục tiêu là tối thiểu hóa tổng chi phí:

\[
\boxed{
C=25\,000x+35\,000y
}
\]
hay:
\[
\boxed{
C=5\,000(5x+7y)\longrightarrow\min.
}
\]

Vì số suất ăn trong thực tế phải là số nguyên nên sau khi xác định miền nghiệm trên mặt phẳng tọa độ, ta sẽ xét các nghiệm nguyên phù hợp.

% =========================================================
% 4.2
% =========================================================

\subsection*{4.2. Biểu diễn các đường biên}

Các đường biên tương ứng với hệ bất phương trình là:
\[
5x+7y=120,
\]
\[
5x+7y=100,
\]
\[
4x+7y=90.
\]
Đưa các phương trình về dạng hàm số theo $x$, ta có:
\[
y=\frac{120-5x}{7},
\]
\[
y=\frac{100-5x}{7},
\]
\[
y=\frac{90-4x}{7}.
\]
Các đường thẳng trên cùng với hai trục tọa độ $Ox$, $Oy$ được sử dụng để xác định miền nghiệm của hệ bất phương trình.

% =========================================================
% 4.3
% =========================================================

\subsection*{4.3. Xác định các đỉnh của miền nghiệm}

Dựa vào các đường biên, ta xác định các giao điểm tạo thành miền nghiệm.
\subsubsection*{Đỉnh A}
Đỉnh $A$ là giao điểm của đường:
\[
5x+7y=100
\]
với trục $Oy$.
Cho $x=0$, ta được:
\[
7y=100
\]
suy ra:
\[
y=\frac{100}{7}.
\]
Vậy:
\[
\boxed{
A\left(0;\frac{100}{7}\right)
}
\approx(0;14,29).
\]
\subsubsection*{Đỉnh B}
Đỉnh $B$ là giao điểm của đường:
\[
5x+7y=120
\]
với trục $Oy$.
Cho $x=0$:
\[
7y=120.
\]
Do đó:
\[
\boxed{
B\left(0;\frac{120}{7}\right)
}
\approx(0;17,14).
\]
\subsubsection*{Đỉnh C}
Đỉnh $C$ là giao điểm của đường:
\[
5x+7y=120
\]
với trục $Ox$.
Cho $y=0$:
\[
5x=120.
\]
Suy ra:
\[
\boxed{C(24;0)}.
\]
\subsubsection*{Đỉnh D}
Đỉnh $D$ là giao điểm của đường:
\[
4x+7y=90
\]
với trục $Ox$.
Cho $y=0$:
\[
4x=90.
\]
Suy ra:
\[
x=\frac{90}{4}=\frac{45}{2}.
\]
Vậy:
\[
\boxed{
D\left(\frac{45}{2};0\right)
}
\approx(22,5;0).
\]
\subsubsection*{Đỉnh E}
Đỉnh $E$ là giao điểm của hai đường:
\[
5x+7y=100
\]
và:
\[
4x+7y=90.
\]
Lấy phương trình thứ nhất trừ phương trình thứ hai:
\[
x=10.
\]
Thay $x=10$ vào phương trình:
\[
5x+7y=100,
\]
ta được:
\[
5(10)+7y=100.
\]
Suy ra:
\[
50+7y=100,
\]
\[
7y=50,
\]
\[
y=\frac{50}{7}.
\]
Do đó:
\[
\boxed{
E\left(10;\frac{50}{7}\right)
}
\approx(10;7,14).
\]
Như vậy, các đỉnh của miền nghiệm là:
\[
\boxed{
\begin{aligned}
A&=\left(0;\frac{100}{7}\right),\\
B&=\left(0;\frac{120}{7}\right),\\
C&=(24;0),\\
D&=\left(\frac{45}{2};0\right),\\
E&=\left(10;\frac{50}{7}\right).
\end{aligned}}
\]

% =========================================================
% 4.4
% =========================================================

\subsection*{4.4. Biểu diễn miền nghiệm trên mặt phẳng tọa độ}

Dựa vào hệ bất phương trình, miền nghiệm được xác định bởi:
\[
\begin{cases}
5x+7y\leq120,\\
5x+7y\geq100,\\
4x+7y\geq90,\\
x\geq0,\\
y\geq0.
\end{cases}
\]
Do đó, miền nghiệm là phần giao của các nửa mặt phẳng tương ứng trong góc phần tư thứ nhất.

Hình dưới đây biểu diễn các đường biên, miền nghiệm và tọa độ các đỉnh của miền nghiệm.

\begin{center}

\begin{tikzpicture}[scale=0.42]

% ---------------------------------------------------------
% Lưới tọa độ
% ---------------------------------------------------------

\draw[step=2,gray!20] (0,0) grid (26,19);

% ---------------------------------------------------------
% Trục tọa độ
% ---------------------------------------------------------

\draw[->,thick] (0,0) -- (27,0)
    node[right] {$x$};

\draw[->,thick] (0,0) -- (0,19)
    node[above] {$y$};

% ---------------------------------------------------------
% Các đường biên
% ---------------------------------------------------------

% 5x + 7y = 120
\draw[blue,thick]
    (0,17.1429) -- (24,0)
    node[pos=0.35,above,sloped]
    {$5x+7y=120$};

% 5x + 7y = 100
\draw[red,thick]
    (0,14.2857) -- (20,0)
    node[pos=0.35,above,sloped]
    {$5x+7y=100$};

% 4x + 7y = 90
\draw[green!60!black,thick]
    (0,12.8571) -- (22.5,0)
    node[pos=0.45,above,sloped]
    {$4x+7y=90$};

% ---------------------------------------------------------
% Tô miền nghiệm
% ---------------------------------------------------------

\fill[orange!25]
    (0,14.2857)
    -- (0,17.1429)
    -- (24,0)
    -- (22.5,0)
    -- (10,7.1429)
    -- cycle;

% ---------------------------------------------------------
% Các đỉnh
% ---------------------------------------------------------

\fill[black] (0,14.2857) circle (0.18);
\fill[black] (0,17.1429) circle (0.18);
\fill[black] (24,0) circle (0.18);
\fill[black] (22.5,0) circle (0.18);
\fill[black] (10,7.1429) circle (0.18);

% ---------------------------------------------------------
% Nhãn các đỉnh
% ---------------------------------------------------------

\node[left] at (0,14.2857)
    {$A\left(0;\frac{100}{7}\right)$};

\node[left] at (0,17.1429)
    {$B\left(0;\frac{120}{7}\right)$};


\node[right] at (24,0)
    {$C(24;0)$};

\node[below] at (22.5,0)
    {$D\left(\frac{45}{2};0\right)$};

\node[right] at (10,7.1429)
    {$E\left(10;\frac{50}{7}\right)$};

% ---------------------------------------------------------
% Tên miền nghiệm
% ---------------------------------------------------------

\node at (15,11)
    {\textbf{Miền nghiệm}};

\end{tikzpicture}

\end{center}

Từ đồ thị, miền nghiệm là một đa giác giới hạn bởi các điểm:

\[
A,\ B,\ C,\ D,\ E.
\]

Các điểm nằm trong miền nghiệm biểu diễn những phương án thỏa mãn đồng thời các điều kiện của bài toán.

% =========================================================
% 4.5
% =========================================================

\subsection*{4.5. Xác định phương án tối ưu}

Hàm mục tiêu của bài toán là:
\[
C=25\,000x+35\,000y.
\]
Ta biến đổi:
\[
C=5\,000(5x+7y).
\]
Theo ràng buộc:

\[
5x+7y\geq100.
\]
Do đó:

\[
C\geq5\,000\times100.
\]
Suy ra:
\[
C\geq500\,000.
\]
Vì vậy, chi phí nhỏ nhất về mặt lý thuyết là:
\[
\boxed{
C_{\min}=500\,000\text{ đồng}.
}
\]
Để đạt được chi phí này, cần:
\[
5x+7y=100.
\]
Đồng thời phải thỏa mãn:
\[
4x+7y\geq90.
\]
Thay:
\[
5x+7y=100
\]
vào điều kiện:
\[
4x+7y\geq90,
\]
ta có:
\[
100-x\geq90.
\]
Suy ra:
\[
x\leq10.
\]
Kết hợp với:
\[
x\geq0,
\]
ta được:
\[
0\leq x\leq10.
\]
Do số suất ăn phải là số nguyên, ta tìm nghiệm nguyên của:
\[
5x+7y=100.
\]
Một nghiệm nguyên phù hợp là:
\[
x=6,\qquad y=10.
\]
Kiểm tra:
\[
5(6)+7(10)=30+70=100,
\]
và:
\[
4(6)+7(10)=24+70=94\geq90.
\]
Vậy:
\[
\boxed{
(x,y)=(6,10)
}
\]
là phương án nguyên phù hợp với mô hình và đạt chi phí tối thiểu.

% =========================================================
% 4.6
% =========================================================

\subsection*{4.6. Kiểm tra phương án tối ưu}

Với: x = 6, y = 10, ta kiểm tra lần lượt các điều kiện.
\subsubsection*{Chi phí}
\[
C=25\,000(6)+35\,000(10).
\]
Do đó:
\[
C=150\,000+350\,000
=500\,000\text{ đồng}.
\]
Vì:
\[
500\,000\leq600\,000,
\]
nên phương án không vượt quá ngân sách.
\subsubsection*{Năng lượng}
Tổng năng lượng là:
\[
600(6)+700(10)
=3\,600+7\,000
=10\,600\text{ kcal}.
\]
Ta có:
\[
10\,600\geq10\,500.
\]
Vì vậy, điều kiện về năng lượng được đáp ứng.
\subsubsection*{Protein}
Tổng lượng protein là:
\[
20(6)+35(10)
=120+350
=470\text{ g}.
\]
Ta có:
\[
470\geq420.
\]
Vì vậy, điều kiện về protein được đáp ứng.
\subsubsection*{Rau xanh}
Tổng lượng rau xanh là:
\[
100(6)+150(10)
=600+1\,500
=2\,100\text{ g}.
\]
Ta có:
\[
2\,100\geq2\,100.
\]
Vì vậy, điều kiện về rau xanh cũng được đáp ứng.

% =========================================================
% 4.7
% =========================================================

\subsection*{4.7. Kết luận kết quả bài toán}

Phương án tối ưu tìm được là:
\[
\boxed{
x=6,\qquad y=10.
}
\]
Nghĩa là trong mô hình giả định, sinh viên lựa chọn:
\begin{itemize}
    \item $6$ suất Thực đơn A;
    \item $10$ suất Thực đơn B.
\end{itemize}
Tổng chi phí là:
\[
\boxed{
500\,000\text{ đồng}.
}
\]
Phương án này đáp ứng đầy đủ các điều kiện:
\[
\boxed{
\begin{aligned}
C&=500\,000\leq600\,000,\\
E&=10\,600\geq10\,500,\\
P&=470\geq420,\\
R&=2\,100\geq2\,100.
\end{aligned}}
\]
Như vậy, mô hình cho phép xác định một phương án thực đơn vừa nằm trong giới hạn ngân sách, vừa đáp ứng các yêu cầu dinh dưỡng đã đặt ra.

Kết quả cũng cho thấy vai trò của hệ bất phương trình hai ẩn trong việc biểu diễn các điều kiện của bài toán và vai trò của phương pháp đồ thị trong việc xác định trực quan miền nghiệm và các điểm biên.

% =========================================================
% KẾT THÚC CHƯƠNG 4
% =========================================================
% =========================================================
% CHƯƠNG 5
% =========================================================

\newpage

\section*{CHƯƠNG 5: ĐÁNH GIÁ VÀ MỞ RỘNG MÔ HÌNH}

\subsection*{5.1. Đánh giá mô hình}

Mô hình toán học được xây dựng trong đề tài đã chuyển một vấn đề
thực tế là lựa chọn thực đơn phù hợp với ngân sách và nhu cầu dinh
dưỡng của sinh viên thành một bài toán quy hoạch tuyến tính với hai
biến.

Thông qua hai biến quyết định, các điều kiện về chi phí và dinh dưỡng
được biểu diễn bằng hệ bất phương trình bậc nhất hai ẩn. Nhờ đó, bài
toán có thể được biểu diễn trực quan trên mặt phẳng tọa độ $Oxy$ và
giải bằng phương pháp hình học.

Kết quả cho thấy phương pháp sử dụng hệ bất phương trình hai ẩn có
thể hỗ trợ xác định những phương án phù hợp với các điều kiện đã đặt
ra. Đồng thời, việc xây dựng hàm mục tiêu giúp lựa chọn phương án có
chi phí thấp hơn trong tập hợp các phương án khả thi.

Mô hình vì vậy phù hợp với mục tiêu của đề tài là minh họa khả năng
ứng dụng kiến thức toán học vào một vấn đề gần gũi với đời sống sinh
viên.

\subsection*{5.2. Ưu điểm của mô hình}

Mô hình có một số ưu điểm nổi bật.
\begin{enumerate}
    \item \textbf{Đơn giản và dễ hiểu:} Mô hình sử dụng hai biến nên
    phù hợp với kiến thức về hệ bất phương trình bậc nhất hai ẩn.
    Bài toán có thể được giải và minh họa trực tiếp trên mặt phẳng
    tọa độ.
    \item \textbf{Có tính trực quan:} Các điều kiện ràng buộc được
    biểu diễn bởi các đường thẳng, từ đó xác định miền nghiệm và các
    điểm biên. Điều này giúp người học dễ dàng quan sát mối quan hệ
    giữa các điều kiện của bài toán.
    \item \textbf{Xem xét đồng thời nhiều yếu tố:} Chi phí, ngân sách
    và các yêu cầu dinh dưỡng được đưa vào cùng một mô hình thay vì
    đánh giá từng yếu tố riêng lẻ.
    \item \textbf{Thể hiện rõ ý nghĩa tối ưu hóa:} Không phải mọi
    phương án đáp ứng điều kiện đều có chi phí giống nhau. Việc xây
    dựng hàm mục tiêu cho phép lựa chọn phương án phù hợp hơn về mặt
    kinh tế.
    \item \textbf{Có khả năng điều chỉnh:} Mô hình có thể được thay
    đổi khi các thông số đầu vào như ngân sách, giá thực phẩm hoặc
    nhu cầu dinh dưỡng thay đổi.
\end{enumerate}

\subsection*{5.3. Hạn chế của mô hình}

Bên cạnh những ưu điểm, mô hình vẫn còn một số hạn chế.

Trước hết, các số liệu về giá thành và giá trị dinh dưỡng của thực đơn
chủ yếu được xây dựng nhằm phục vụ mục đích học tập và minh họa.
Trong thực tế, giá thực phẩm có thể thay đổi theo thời gian, địa điểm
mua hàng và chất lượng sản phẩm.

Thứ hai, mô hình chỉ xét một số tiêu chí dinh dưỡng cơ bản. Một thực
đơn thực tế còn cần quan tâm đến nhiều yếu tố khác như chất béo,
chất xơ, vitamin, khoáng chất và sự cân đối giữa các nhóm thực phẩm.

Thứ ba, mô hình sử dụng hai biến nên số lượng loại thực phẩm được xét
còn hạn chế. Trong thực tế, sinh viên có thể lựa chọn nhiều loại thực
phẩm và món ăn khác nhau.

Ngoài ra, mô hình chưa xét đến các yếu tố cá nhân như sở thích ăn
uống, dị ứng thực phẩm, thói quen sinh hoạt hoặc nhu cầu dinh dưỡng
riêng của từng người.

Vì vậy, kết quả của mô hình chỉ có ý nghĩa trong phạm vi các giả định
được đặt ra và không nên được xem là một thực đơn dinh dưỡng áp dụng
chung cho mọi sinh viên.

\subsection*{5.4. Ý nghĩa thực tiễn của mô hình}

Mặc dù còn đơn giản, mô hình cho thấy hệ bất phương trình hai ẩn có
thể được sử dụng để hỗ trợ việc lập kế hoạch chi tiêu trong đời sống.

Đối với sinh viên có ngân sách giới hạn, việc lựa chọn thực phẩm
không chỉ phụ thuộc vào giá thành mà còn cần xem xét giá trị dinh
dưỡng. Một phương án có giá thấp chưa chắc đã đáp ứng được yêu cầu
về năng lượng và protein. Ngược lại, một phương án có nhiều chất dinh
dưỡng hơn có thể dẫn đến chi phí cao hơn.

Việc mô hình hóa giúp biểu diễn đồng thời các yêu cầu này bằng ngôn
ngữ toán học. Nhờ đó, người sử dụng có thể xác định những phương án
phù hợp và so sánh chi phí giữa các phương án khác nhau.

Qua đó, đề tài cho thấy toán học có thể đóng vai trò hỗ trợ trong quá
trình lập kế hoạch, phân bổ ngân sách và đưa ra quyết định trong
những vấn đề thực tế.

\subsection*{5.5. Tính linh hoạt của mô hình}

Một ưu điểm quan trọng của mô hình là có thể thay đổi các thông số
mà không cần xây dựng lại toàn bộ phương pháp giải.

Chẳng hạn, khi ngân sách của sinh viên thay đổi, bất phương trình
biểu diễn ngân sách có thể được điều chỉnh. Nếu ngân sách tăng, miền
nghiệm có thể mở rộng và xuất hiện thêm nhiều phương án lựa chọn.
Ngược lại, nếu ngân sách giảm, miền nghiệm có thể thu hẹp.

Tương tự, khi nhu cầu dinh dưỡng thay đổi, các ràng buộc về năng
lượng hoặc protein cũng có thể được điều chỉnh. Khi đó, vị trí các
đường biên thay đổi và miền nghiệm có thể thay đổi theo.

Giá thực phẩm cũng là một yếu tố có thể thay đổi. Khi giá của một loại
thực đơn tăng hoặc giảm, hệ số trong hàm mục tiêu sẽ thay đổi. Điều
này có thể làm cho phương án tối ưu thay đổi mặc dù các điều kiện
dinh dưỡng vẫn giữ nguyên.

Như vậy, mô hình có tính linh hoạt và có thể được sử dụng để phân
tích nhiều tình huống khác nhau liên quan đến ngân sách và nhu cầu
dinh dưỡng.

\subsection*{5.6. Hướng mở rộng mô hình}

Để mô hình gần với thực tế hơn, có thể phát triển theo một số hướng.

Trước hết, có thể tăng số lượng loại thực phẩm hoặc thực đơn được đưa
vào mô hình. Khi đó, thay vì chỉ sử dụng hai biến $x$ và $y$, có thể
sử dụng các biến $x_1,x_2,\ldots,x_n$ để biểu diễn số lượng của nhiều
loại thực phẩm khác nhau.

Ví dụ, có thể đưa vào mô hình các nhóm thực phẩm như cơm, thịt, cá,
trứng, sữa, rau xanh và trái cây. Khi số lượng biến tăng lên, bài toán
sẽ trở thành một bài toán quy hoạch tuyến tính nhiều biến và có thể
được giải bằng các công cụ tính toán phù hợp.

Bên cạnh đó, có thể bổ sung thêm các ràng buộc dinh dưỡng như chất
béo, chất xơ, vitamin và khoáng chất. Điều này giúp mô hình phản ánh
đầy đủ hơn các yêu cầu của một thực đơn cân đối.

Một hướng phát triển khác là sử dụng dữ liệu thực tế về giá thực phẩm
tại một khu vực cụ thể. Khi đó, mô hình có thể được cập nhật theo giá
thị trường và có khả năng phản ánh tốt hơn điều kiện chi tiêu thực tế
của sinh viên.

Ngoài ra, có thể bổ sung các ràng buộc về số lượng suất ăn hoặc tần
suất sử dụng một loại thực phẩm. Chẳng hạn, có thể quy định số suất
ăn trong tuần phải đạt một mức nhất định hoặc giới hạn số lần sử dụng
một loại thực đơn.

Cuối cùng, mô hình có thể được mở rộng theo hướng xem xét nhiều mục
tiêu cùng lúc, chẳng hạn vừa tối thiểu hóa chi phí vừa đảm bảo mức
dinh dưỡng phù hợp. Tuy nhiên, khi đó bài toán sẽ phức tạp hơn và cần
sử dụng những phương pháp tối ưu hóa nâng cao.


% =========================================================
% KẾT THÚC CHƯƠNG 5
% =========================================================
%=========================================================
% TÀI LIỆU THAM KHẢO
%=========================================================

\clearpage
\addcontentsline{toc}{chapter}{TÀI LIỆU THAM KHẢO}

\begin{thebibliography}{99}

\bibitem{openstax}
OpenStax (2023),
\textit{Intermediate Algebra 2e, Section 3.4: Graph Linear Inequalities in Two Variables},
OpenStax, 
\url{https://openstax.org/books/intermediate-algebra-2e/pages/3-4-graph-linear-inequalities-in-two-variables}.

\bibitem{vietyen}
Trường THPT Việt Yên số 2, Bắc Ninh,
\textit{Chuyên đề ứng dụng hệ bất phương trình bậc nhất hai ẩn vào giải một số bài toán thực tế},
Trường THPT Việt Yên số 2, Bắc Ninh,
\url{https://thptvietyenso2.bacninh.edu.vn/tai-nguyen/tai-lieu/chuyen-de-ung-dung-he-bat-phuong-trinh-bac-nhat-hai-an-vao-giai-mot-so-bai-toan-thuc-te-..html}.

\bibitem{huynhngai}
Huỳnh Văn Ngãi, Nguyễn Văn Vũ,
\textit{Lý thuyết tối ưu},
QNU,
\url{https://ahust.vn/tai-lieu/xem/ly-thuyet-toi-uu-huynh-van-ngai-nguyen-van-vu-qnu}.

\bibitem{nutrihome}
Nutrihome,
\textit{Khẩu phần ăn hợp lý},
Nutrihome,
\url{https://nutrihome.vn/khau-phan-an-hop-ly/}.

\bibitem{toan10}
Bộ Giáo dục và Đào tạo,
\textit{Sách giáo khoa Toán 10},
Nhà xuất bản Giáo dục Việt Nam.

\bibitem{chatgpt}
OpenAI,
\textit{ChatGPT},
Công cụ trí tuệ nhân tạo hỗ trợ nghiên cứu và xây dựng nội dung,
\url{https://chatgpt.com/}.

\bibitem{gemini}
Google,
\textit{Gemini},
Công cụ trí tuệ nhân tạo hỗ trợ nghiên cứu và tham khảo thông tin,
\url{https://gemini.google.com/}.

\bibitem{googleai}
Google,
\textit{Google AI},
Công cụ trí tuệ nhân tạo hỗ trợ nghiên cứu,
\url{https://ai.google/}.

\end{thebibliography}
\end{document}








































































































































































































































